\par\noindent\textbf{(1)} 取ψ=ψ₀e\textasciicircum{}\{i(kx+ly−ωt)\}，有 ∇²ψ=−(k²+l²)ψ。代入得 iω(k²+l²+L\_D\textasciicircum{}\{-2\})ψ+iβkψ=0，故 ω=−βk/(k²+l²+L\_D\textasciicircum{}\{-2\})。β>0、k>0时频率为负，表现西向相传播。
\par\noindent\textbf{(2)} x向相速度 c\_px=ω/k=−β/(k²+l²+L\_D\textasciicircum{}\{-2\})<0。二维群速度分量为 c\_gx=∂ω/∂k=β(k²−l²−L\_D\textasciicircum{}\{-2\})/(k²+l²+L\_D\textasciicircum{}\{-2\})²、c\_gy=∂ω/∂l=2βkl/(k²+l²+L\_D\textasciicircum{}\{-2\})²；群速可与相速东西向相反。
\par\noindent\textbf{(3)} 由于 c\_gx 的符号等于k²−l²−L\_D\textasciicircum{}\{-2\}，短x波满足 k²>l²+L\_D\textasciicircum{}\{-2\} 时群能量向东传播，长x波时向西。零群速条件是 k²=l²+L\_D\textasciicircum{}\{-2\}；不能仅据相位总向西就断言能量总向西。
\par\noindent\textbf{(4)} 当L\_D→∞时回复正压Rossby色散ω=−βk/(k²+l²)，k,l→0存在非一致极限；若固定L\_D有限且k→0，ω≈−βk/(l²+L\_D\textasciicircum{}\{-2\})。β→0则频率趋零，表明Rossby恢复机制来自行星涡度梯度。
